% !TeX root = ../main.tex

\xchapter{绪论}{Introduction}

\xsection{研究背景与意义}{Research Background and Significance}

数学教育是基础教育体系的核心环节，其中平面几何课程在培养学生逻辑推演与结构认知能力方面发挥着不可替代的作用。在实际的几何解题过程中，学生需要观察图形结构、抽象几何约束条件，并依托公理体系完成严密的逻辑演绎。这种视觉感知与形式化推演相融合的思维训练，能够有效促进学生科学探究能力的形成。近年来，随着国家教育数字化转型战略的持续推进，如何借助信息技术赋能数学教育，已成为基础教育领域的核心研究议题。在此背景下，平面几何问题自动求解逐渐成为检验智能系统认知能力的关键标尺，在智能辅导、自动解题等教育数字化应用场景中受到了广泛关注。如图\ref{fig:problem_example}所示，一个典型的平面几何问题求解任务通常包含几何图像、自然语言描述的已知条件及待求目标，这要求智能系统具备强大的跨模态理解与逻辑推理能力。

随着人工智能技术的快速发展，几何问题求解方向的研究已取得显著进展。在数据建设层面，几何数据集完成了从早期纯文本题库向多模态架构的跨越，部分前沿基准已引入程序化描述语言与分步式结构化标注体系。在问题求解范式层面，底层算法同样完成了快速的技术迭代：从早期依赖人工定义规则、通过符号引擎执行严密逻辑证明的传统范式，逐步转向基于端到端神经网络实现多模态图文信息联合感知的智能范式。近年来，大语言模型凭借其优异的上下文语义解析与泛化能力，大幅突破了传统算法的推理性能上限。值得注意的是，大语言模型的突破性进展，也反向推动了几何数据基准的演进升级：研究人员开始利用大模型强大的生成与解析能力，自动化扩充测试样本、解构复杂推理步骤；而这些蕴含丰富逻辑链条的垂直领域数据，又被进一步用于夯实大模型在几何场景下的认知能力，形成了数据与模型协同演进的良性生态。

\begin{figure}[htbp]
    \centering
    \includegraphics[width=0.9\textwidth]{c1_example.png}
    \caption{平面几何问题求解任务样例}
    \label{fig:problem_example}
\end{figure}


尽管求解技术与数据基准在相互促进中持续发展，但当前主流的大语言模型在应对长推理链条的复杂几何问题时，仍存在一定的局限性。具体而言，当面对涉及多步演绎的综合题型时，模型在复杂的跨模态场景下，容易出现图像理解失效或图文语义不匹配等问题。这些早期感知层面的微小偏差不仅难以被模型自我修复，还会由于推理步骤间的前后依赖关系而被不断放大，最终导致逻辑链条偏离并引发推理错误。如图\ref{fig:llm_error}所示，即使是较为先进的OpenAI o4-mini~\cite{openai2025o3o4mini}模型，在进行长链条几何推理时仍然会发生推理失效问题。

\begin{figure}[htbp]
    \centering
    \includegraphics[width=0.8\textwidth]{c1_llm_error.pdf}
    \caption{大模型在长步骤平面几何题目中的推理失效示例}
    \label{fig:llm_error}
\end{figure}

深入剖析该案例可以发现，由于当前大模型在推理时缺乏对中间步骤的有效控制或校验机制，早期产生的微小感知错误（如步骤2）会被直接代入后续的长链条计算环节中。在这种情况下，初始的误差会随着推演步骤的增加而被不断累积，最终导致求解结果发生偏离。这种因中间过程难以校验而引发的误差累积问题，直接制约了智能解题工具在真实教育场景中的可靠落地。

为突破上述技术瓶颈，本文聚焦长链条几何问题自动求解这一核心痛点，开展了系统性的数据基准构建与算法设计研究。一方面，本文构建了贴近真实教育场景的高难度多模态几何评测基准，补充了细粒度的推理步骤与形式化语义标注；另一方面，提出了一种融合神经感知与弱符号规则反馈的两阶段增强方法，通过引入过程级校验机制，有效抑制了长链推演中的误差累积与逻辑发散。在此基础上，本文的研究不仅在理论层面探索了可验证的神经符号协同新机制，验证了中间过程约束在缓解多模态大模型逻辑幻觉、提升长程推理稳定性方面的关键作用；更在实际应用层面弥补了复杂长链条推理评测资源的空白，为构建高可靠、可解释的智能辅导系统提供了坚实底座，有力助力自动解题与个性化导学等数字教育工具在真实数学课堂中的落地应用。

\xsection{国内外研究现状}{Related Work}

近年来，几何问题求解作为多模态理解与复杂逻辑推理交叉的核心领域，受到了学术界的广泛关注。为了更清晰地呈现该领域的发展历程，本节将从“数据基准”和“求解方法”两个主要维度对现有工作展开介绍。
% 具体的分类体系与代表性研究如图~\ref{fig:related_work} 所示。
下文首先介绍几何数据集从早期人工标注向大模型自动化合成的发展过程，以及几何文本与图像的常见解析方法，随后将对比基于规则、神经网络以及大语言模型的三大类主流几何问题求解方法。

% \begin{figure}[htbp]
%     \centering
%     \includegraphics[width=1.0\textwidth]{c1_related_work.png}
%     \caption{几何问题求解研究相关工作}
%     \label{fig:related_work}
% \end{figure}

\xsubsection{几何问题数据基准}{Benchmarks for Geometric Problems}

随着多模态推理技术的发展，几何评测基准的构建方式经历了从耗时的人工收集到借助大语言模型实现自动化收集、标注的演变。当前几何数据集的构建主要可划分为人工构造、大模型辅助标注、大模型辅助增强和大模型辅助合成四种范式。

\subsubsection{人工构造}

早期的几何数据集主要依赖领域专家手动构建，通过整合教材等来源的题目并进行精细化标注，确保了数据的高准确性和对教学场景的贴近性。然而，这种高度依赖手工介入的范式也导致数据集规模普遍较小，且扩展成本较高。例如，经典数据集Geometry3K~\cite{lu2021intergps} 从6至12年级的教材中整理了3002道SAT级别的几何题，并为每道题提供了详细的形式化逻辑标注，为神经符号系统中的自动化推理奠定了基础；GeoQA~\cite{chen2021geoqa} 则包含4998道图文结合的几何题，不仅通过人工方式爬取了大量带图题目，还为每道题编写了可执行的求解程序，推动了早期多模态理解与数值推理的研究；随着对深层逻辑推理能力评测的需求增加，单纯的数值计算已无法满足评测需要，因此 UniGeo~\cite{chen2022unigeo} 引入了9543道带有完整推导步骤的证明题，首次在一个统一的序列生成框架下同时支持几何计算和逻辑证明两类任务；为了减少自然语言描述带来的歧义，FormalGeo7k~\cite{zhang2023formalgeo} 基于严格的几何形式化体系，为近7000道题的文本条件、图形元素以及定理使用顺序，人工编写了结构化的形式化语言表示，从而实现对机器证明过程的精确验证；此外，为了更全面地评估当前的大语言模型和多模态模型，GeoEval~\cite{zhang2024geoeval} 整合了7个主流公开数据集，并加入人工设计的高难度题目，构建了一个覆盖平面几何、立体几何和解析几何，同时包含反向推理和挑战性变体的综合性评测基准。

\subsubsection{大模型辅助标注}

为了突破人工结构化标注的成本瓶颈，研究者开始引入大模型辅助标注范式。该方法利用大语言模型对现有基础题目进行深度理解，自动生成细粒度的中间推理步骤，从而将简单的解答转化为高质量的结构化数据。例如，GeoSense~\cite{xu2025geosense} 提出了一种结合 GPT-4o 与专家人工校验的半自动化标注流程，从而构建了包含5层架构的几何知识体系的评测基准；GPSM4K~\cite{jaiswal2024gpsm4k} 借助多模态大模型完成原始几何题目的变体生成，以及对解题过程进行了模块化重构；GNS260K~\cite{ning2025gns} 利用 OpenAI GPT-4，以符号化解析结果和符号解法程序为输入，生成了自然语言形式的推理轨迹，从而将解答中的关键思维节点转化为高质量的 CoT 标注。这标志着几何推理数据构建正从依赖高成本的人工专家标注，转向利用大模型实现自动化、规模化生产的新范式。

\subsubsection{大模型辅助增强}

在辅助标注的基础上，大模型辅助增强范式进一步拓宽了数据构建的思路。该范式通常以少量原始题目为“种子”，借助大模型进行条件变换、提问方式改写或数据清洗等变异操作，从而低成本地横向拓展数据集的规模与多样性。例如，G-LLaVA~\cite{gao2023gllava} 团队基于现有的少量几何数据，利用大语言模型进行图文对齐生成与增强，构建了规模达17万的 Geo170K 训练数据，极大提升了基准对几何特征的覆盖面；MathVerse~\cite{zhang2024mathverse} 提出多条件增强，利用大模型将 2612 道题增强为 15000+ 道具有不同“信息泄露”程度的变体样本；MathV360K~\cite{shi2024mathllava} 利用GPT-4V对原始种子数据进行图像挖掘生成新问答题，同时还生成复杂化题目以提升推理深度。

\subsubsection{大模型辅助合成}

作为目前最前沿的构建方式，大模型辅助合成标志着数据生产走向了全链路自动化。它完全脱离了对现有人工题库的依赖，通过结合符号引擎与大语言模型的生成能力，能够从零开始自动化合成海量的“条件-图形-解答”元组，甚至直接将模型的探索试错过程转化为数据集。其中最具代表性的是 AlphaGeometry~\cite{trinh2024alphageometry}，该工作通过神经符号系统合成上亿条不同复杂度的定理和证明路径，首次实现了无需人类示范的奥林匹克级别高质量几何数据合成。GeoDANO~\cite{cho2025geodano} 基于 AlphaGeometry 提出了一种基于程序化合成的数据引擎，利用形式化语言与渲染系统生成了包含 20 万对图表与描述的预训练数据集，同时将真实场景数据转化为合成风格。MathCoder-VL~\cite{wang2025mathcodervl} 则进一步将这一范式推向了“全栈自动化”的高度。该工作不仅利用 GPT-4o 进行数据清洗与模态翻译，更核心的是利用 Qwen2.5~\cite{qwen25} 系列大模型构建了合成数据引擎，通过以高随机性参数驱动图像生成与题目重写，实现了从“旧图”到“新题”的创造性改造，打破了人工题库的规模瓶颈。

% ---------------------------------------------------------------
\xsubsection{几何文本与图像解析}{Geometric Text and Diagram Parsing}

几何问题的多模态特性决定了其解析任务包含图像侧的感知与文本侧的语义理解两个核心维度。本小节主要介绍如何将非结构化的几何题目转化为计算机可处理的结构化符号体系或神经网络特征表示。

\subsubsection{几何文本解析与表示}
文本语义解析的任务是将自然语言描述的几何条件映射为严格的形式化谓词或密集的语义表示，进而输入符号引擎、神经求解器或大语言模型，以实现几何问题的自动化推理。传统方法主要基于关键词匹配或预定义的语法规则，难以应对自然语言的多样性与歧义性。Inter-GPS~\cite{lu2021intergps} 率先引入形式化语言作为符号推理引擎的输入表示，并通过正则表达式匹配实现语义转换；GeoDRL~\cite{peng2023geodrl} 等后续工作亦沿用了这一基于规则的解析范式。AlphaGeometry~\cite{trinh2024alphageometry} 则通过生成海量的纯形式化合成数据，训练语言模型直接在严格的逻辑符号空间内进行几何构造与推理；其核心推理引擎并不直接处理自然语言，在实际应用中需依赖外部模块或人工将自然语言预先转化为指定的领域形式语言。

随着深度学习的发展，神经机器翻译范式被引入该领域，如 UniGeo~\cite{chen2022unigeo} 等工作将文本语义解析建模为序列到序列（Seq2Seq）的生成任务，直接将自然语言转化为可执行的数学表达式或程序序列。另一方面，在基于密集特征的端到端神经方法中，系统将自然语言条件输入至预训练文本编码器中，直接提取出高维的连续语义向量，以便进行深度的多模态特征融合。例如，GeoQA~\cite{chen2021geoqa} 利用预训练语言模型提取问题文本的深层语义特征，并通过共同注意力（Co-attention）机制将其与图像的卷积特征进行动态交互与对齐；PGPSNet~\cite{zhang2023pgps} 则在此基础上，将自然语言问题与从配图中解析出的结构化文本子句相拼接，共同输入至 RoBERTa 模型中以获取增强的联合语义表示，从而为后续的跨模态推理提供了更完备的线索。

\subsubsection{几何图像解析与特征提取}

几何图像解析旨在从题目配图中提取基础图元及其拓扑关系，并将其转化为可供推理计算的模型输入。对于需要显式识别几何基元的工作来说，对图中符号与关键点标签的高精度识别是实现后续图文逻辑对齐的先决条件，这通常高度依赖于以 PaddleOCR~\cite{cui2025paddleocr} 为代表的光学字符识别（OCR）技术。在早期研究中，获取图形形式化表达的途径多依赖于外部的显式解析模块，如早期符号推理工作 Inter-GPS~\cite{lu2021intergps}，这类方法在面对复杂图像背景时往往缺乏鲁棒性。为了提升底层感知的端到端能力，PGDPNet~\cite{zhang2022pgdp} 提出了一种神经网络架构，能够直接从图片中预测并解析出离散的图元坐标与几何约束，输出形式化语言以及几何基元关系，初步确立了基于神经网络提取图形结构信息的基础。

随着多模态推理技术的发展，当前的主流图像解析范式已全面转向深度神经网络，致力于将图形转化为深层的连续特征向量，以实现更加高效的跨模态交互。例如，SCA-GPS~\cite{ning2023scagps} 专门引入了图形编码器，通过将图像划分为局部图块来提取视觉特征，并在保留局部布局线索的同时与文本特征进行拼接对齐；LANS~\cite{li2024lans} 则设计了布局感知的神经架构，在特征融合层面隐式地捕获了图元间的空间几何分布。为了进一步结合显式拓扑先验与隐式神经表征的优势，最新的研究开始引入图神经网络对几何配图进行结构化编码。例如，$\mathrm{G^3PS}$~\cite{wu2026geograph} 创新性地将几何配图转化为几何图（Geo-Graph），把解析出的图元作为节点，融合其语义与空间位置信息，随后利用强大的图 Transformer等神经机制对图元间的复杂拓扑关联进行连续编码。这种利用神经网络进行向量化特征提取与深层解析的演进，使得图形的全局拓扑与局部细节能够在隐空间中得到完整保留，为后续的深度推理提供了丰富的语义线索。

% \subsubsection{多模态表示融合与对齐}
% 多模态信息的解析与融合主要分为两类技术路线：结构化符号表示对齐与神经网络特征融合。 在结构化表示层面，目标是将图像侧与文本侧的形式化结果整合至统一的谓词逻辑空间，形成完整的问题约束集合。除了 Inter-GPS~\cite{lu2021intergps} 所采用的谓词规范外，E-GPS~\cite{E-GPS: Explainable Geometry Problem Solving via Top-Down Solver and Bottom-Up Generator} 构建的图文统一形式化表征使得求解引擎可以直接调用多模态互补条件。在最新进展中，Causal-R~\cite{Causal-R: A Causal-Reasoning Geometry Problem Solver for Optimized Solution Exploration} (NeurIPS 2025) 将对齐后的多模态符号图转化为因果推理图，通过矩阵推演压缩搜索空间，大幅提升了求解效率。在神经网络特征表示层面，研究者通常采用经典的编码器-解码器（Encoder-Decoder）架构来实现多模态对齐。系统利用编码器分别提取图文特征后，通过跨模态注意力等融合机制进行深度语义对齐，随后交由解码器直接映射到最终答案空间或中间形式化语言（如 PGPSNet~\cite{A multi-modal neural geometric solver with textual clauses parsed from diagram} 和 DFE-GPS~\cite{Diagram formalization enhanced multi-modal geometry problem solver}）。然而，不论是结构化解析还是端到端神经融合，当前技术在处理隐性空间拓扑关系（如点在线段上的相对顺序）以及极端视觉干扰下仍面临挑战，探索更具强健性的图文双向校验机制依然是缓解长链条推理逻辑失稳的关键。


% ---------------------------------------------------------------
\xsubsection{几何问题求解方法}{Methods for Geometric Problem Solving}

在求解算法方面，现有的几何推理研究成果可以归结为三大主要技术路线，即基于规则的方法、基于神经网络的方法以及基于大语言模型的方法。它们在逻辑严谨性与特征泛化能力之间有着不同的权衡与侧重。

\subsubsection{基于规则的方法}

早期平面几何问题求解主要遵循符号推理路线，即基于规则的方法。这类方法的核心思想是将几何解题转化为基于形式化语言的严格演绎过程。其系统通常具备极高的确定性与透明度，高度依赖预先设定的几何公理、定理体系以及专家手工制定的搜索规则，主要通过符号计算引擎执行逻辑推演。经典工作 Inter-GPS~\cite{lu2021intergps} 依靠显式的形式化语言解析模块和外置的前向搜索定理库来执行推演；E-GPS~\cite{wu2024egps} 创新性地提出了自顶向下的求解器，通过从目标出发反向推理并结合自底向上的问题生成器来增强定理预测能力；GeoTree~\cite{wang2026geotree} 构建了基于蒙特卡洛树搜索（MCTS）的动态推理框架，通过将“定理预测”与“定理应用”解耦，利用大语言模型进行定理探索，同时结合符号求解器进行严格的逻辑推演。尽管此类方法在求解过程中具有完美的解释性和绝对的逻辑可靠性，但它依赖于严格的形式化输入，难以直接处理非结构化的自然语言或复杂的图像像素输入，且其泛化应用能力严重受限于手工规则的完备性。

\subsubsection{基于神经网络的方法}

随着深度学习技术的兴起，几何求解范式经历了从“符号演绎”向“特征拟合”的转变，基于神经网络的方法逐渐受到广泛关注。此类方法通常采用经典的编码器-解码器（Encoder-Decoder）架构，旨在实现端到端的多模态自动求解。具体而言，系统通常利用强大的神经网络编码器（如CNN或Transformer）提取几何图像的视觉特征与文本的语义表示，进行跨模态特征融合；随后交由解码器直接映射到最终的答案空间或中间形式化语言，从而跨过了早期繁琐的人工规则解析阶段。例如，PGPSNet~\cite{zhang2023pgps}利用深度神经网络特征融合模块，将视觉图元特征与文本编码直接映射到最终的答案空间；DFE-GPS~\cite{zhang2025dfegps} 引入了“图形式化”中间层，通过构建几何形式语言作为视觉与语义的桥梁，利用合成数据增强视觉编码器对几何结构的理解；NS-GPS~\cite{wang2025reasoning} 则构建了一种神经符号协同框架，通过将神经网络的“感知能力”与符号计算的“精确逻辑”解耦并有机结合，在利用神经网络生成子程序的同时，依赖外部符号执行器进行严格的数值计算与中间反馈。然而，由于缺乏显式的逻辑路径约束，这些基于纯神经网络的模型在面对长步骤、高复杂度的证明题时，往往表现出“黑盒”特性，难以保证推理过程的严谨性与连贯性。

\subsubsection{基于大语言模型的方法}

近年来，大语言模型及多模态大模型的爆发，为几何问题求解带来了革命性的范式跃迁。与传统深度学习模型依赖特定任务数据的端到端拟合不同，大模型凭借海量预训练数据中涌现出的强大上下文理解与通用逻辑推理能力，能够以更加拟人化的思维链（Chain-of-Thought）进行多步解题。当前的研究生态呈现出多元化发展趋势，既包含探究如何通过提示工程直接激发通用多模态大模型零样本解题潜能的尝试，也催生了一批针对数学与几何垂直领域深度定制的专有大模型。相关研究主要沿着两大路径演进：一是推理范式创新，侧重于优化模型的逻辑推导与解答过程；二是模型微调与预训练，侧重于从底层架构为模型注入特定的几何专业能力。

侧重于推理范式创新的方法，往往通过特定的能力微调结合复杂的提示策略，来优化模型的逻辑推导与上下文生成过程。 例如，Geo-LLaVA~\cite{xu2024geollava} 在基础训练之上提出了元上下文学习策略，通过在提示中输入特定的几何图文配对示例，有效激发了大模型对空间关系和几何求解的直观联想能力；GF-Reasoner~\cite{yang2025GFReasoner} 框架经由后训练（Post-training） 将自然语言思维链与形式化几何代码相结合，使大模型的推导过程兼具解释性与计算严谨性；GOLD~\cite{zhang2024gold} 则将几何图元转化为自然语言描述，证明了自然语言形式的思维链能显著提升大模型的对齐和解题能力。

基于监督微调与预训练的方法也逐渐成为几何推理领域的主流增强范式，通过更新模型参数，从底层架构上为通用大模型注入专业的几何解题能力。例如，EAGLE~\cite{li2024eagle} 提出了由大语言模型赋能的视觉指令微调框架，有效提升了多模态模型对“点、线、角”等基础几何元素的细粒度感知能力；GNS~\cite{ning2025gns} 则利用指令微调使模型学会输出结构化的“神经符号轨迹”，实现了视觉感知与符号逻辑的深度耦合。预训练层面，GeoX~\cite{xia2024geox} 引入了专门针对几何领域的统一形式化图文预训练范式，针对性地训练了图表编码器和符号解码器，为大模型提供几何先验知识。

\xsection{研究挑战}{Research Challenges}

平面几何问题求解是检验人工智能多模态理解与逻辑推理能力的重要标准。当前的研究正从简单的浅层计算向复杂的多步推理发展。在实际的数学教育和研究中，几何问题的证明通常需要多次中间结论转换和连续的定理调用，这种“长链条”推理不仅更符合人类的数学思维模式，也对模型在长程逻辑推理的连续性和复杂约束处理能力提出了更高要求。基于当前多模态理解和神经符号推理的研究现状，以及平面几何对上下文逻辑的高度依赖特性，本文总结出模型在长链条推理中面临的两个核心挑战：

第一，长链条推理的过程评测缺失与图文形式化瓶颈。要客观评价模型是否具备真正的长链条推理能力，不能仅看最终答案，必须深入核验其每一个中间推导步骤。尽管当前几何数据集的规模在不断扩大，但由于缺乏高难度综合题的细粒度过程标注，现有的算法评估很难深入到具体的推理步骤层面。同时，真实的几何题目包含简洁的自然语言和直观图形，将这些非结构化的文本和蕴含隐含拓扑关系的图像，无损且精准地转化为计算机能够自动校验的严格形式化逻辑，存在较高的技术门槛。自动化转换链路的尚不完善，成为了连接自然语言与符号逻辑的瓶颈。这种形式化转换桥梁与标准化过程评测平台的匮乏，直接制约了对长链条推理中间过程的科学评测，也限制了推理过程自动校验以及混合智能求解模型的后续研究。

第二，跨模态长链条推理中的逻辑失稳与误差累积风险。平面几何的解题过程具有极强的前后依赖性。在长达数十步的推演中，早期图文理解阶段产生的哪怕极其微小的感知偏差，都会随着推导深度的增加而被迅速放大和累积，极易导致全盘逻辑崩溃。此外，当模型试图同时处理抽象的逻辑推理和复杂的视觉图像时，两者容易产生耦合干扰：视觉信息的引入有时反而会干扰模型学习复杂的逻辑推演关系，使其很容易受到图中冗余线条或非精确标注的干扰。现有的端到端生成模型大多缺乏针对中间步骤的动态校验机制，在没有外部即时反馈的情况下，模型一旦出错就容易沿着错误方向继续推演，难以硬性阻断误差的向下发散，从而导致模型难以在复杂的跨模态输入下建立并维持稳定的长链条演绎能力。

\xsection{研究内容}{Research Contents}

本论文工作依托于教育部基础学科和交叉学科突破计划(JYB2025XDXM116)及国家自然科学基金(62477036)，主要关注于几何问题的长链条推理能力的评估与增强，是相关研究计划的核心子任务。本文通过对几何问题多模态信息的深度解构，将非结构化的几何图文转化为结构化的符号推理序列，并构建起一套面向长链条推理、具备细粒度逻辑标注的多模态评测基准；在此基础上，引入结合弱符号反馈的两阶段强化训练算法，实现了对复杂证明路径的精准约束与推理轨迹的有效引导。研究成果构建了针对几何长链条推理的增强范式，为智慧教育提供了核心逻辑支撑，助力几何教学实现深度智能化变革。整体研究框架如图\ref{fig:thesis_arch}所示：

\begin{figure}[htbp]
    \centering
    \includegraphics[width=1.0\textwidth]{c1_thesis_arch.pdf}
    \caption{研究内容框架}
    \label{fig:thesis_arch}
\end{figure}


\subsubsection{面向长链条推理的多模态平面几何基准构建}
针对长链条推理的过程评测缺失与图文形式化瓶颈，本文构建了面向长步骤平面几何推理的多模态评测基准 GeoLStep。本研究从真实的中考试卷中系统收集并筛选了2097道具有较高难度的平面几何综合题，涵盖了数值计算与结论证明两类核心任务。在数据处理方面，借助多模态大语言模型与光学字符识别技术，对题目图文和标准解析进行了高精度的内容校对与细粒度逻辑解构，并补充了完整的推理步骤数以及辅助线依赖等结构化标签。为突破图文形式化瓶颈，本文系统性地开发了自然语言与几何图像通向严格符号逻辑的形式化转换链路：一方面利用指令微调的小参数大语言模型实现文本到谓词的高效转换，另一方面结合微调后的视觉检测模块与图元解析网络实现图像特征的符号解析。在此基础上，针对过程评测缺失的问题，本工作还引入了基于大模型的自动化过程评估机制，确立了答案准确率与过程逻辑正确率双重评价标准。该基准不仅弥补了当前长链条逻辑评测数据的空白，也为后续复杂过程级别的模型对比与强化训练提供了坚实的数据底座。

\subsubsection{面向长链条平面几何推理的两阶段增强方法}
针对跨模态长链条推理中的逻辑失稳与误差累积风险，本文提出了面向长链条几何问题求解的两阶段增强方法，并设计了一套结合组相对策略优化（GRPO）与弱符号验证器的强化学习框架。在第一个训练阶段进行文本推理增强，模型在剥离视觉要素干扰的纯文本环境中，利用外部工具链生成的结构化谓词进行训练，集中学习基础的演绎范式与步骤格式规范，从而建立稳定的长链条推理先验，有效缓解视觉特征耦合带来的逻辑失稳。在第二个阶段进行多模态多粒度过程强化，模型恢复处理多模态图文联合输入。为硬性阻断长程推演中的误差级联放大，系统同步引入包含解析器与谓词融合器的弱符号验证器，对模型自回归生成的每一步中间结论进行即时的合法性检验与冗余判定。通过将上述步骤级反馈转化为策略优化的局部奖励信号，并施加参考策略的分布约束，该方法能够有效规范大模型的生成轨迹，显著抑制了多步推演后期的误差发散现象，大幅提升了模型在复杂多模态输入下的推理连续性与求解准确率。

\subsubsection{平面几何问题求解系统开发}
为了整合落地平面几何问题求解相关工作，本文开发并实现了交互式的平面几何问题求解原型系统。本系统使用基于前后端解耦的分层式工程架构，将前序构建的数据集资源与核心强化求解算法进行了完整的系统级集成。平台前端支持用户便捷地上传“题干加配图”的多模态测试数据；后端则采用基于虚拟环境隔离的部署方案，允许并发调度传统树状搜索、演绎可解释方法以及微调大语言模型等多个不同范式的求解器。系统特别设计了异步通信与分阶段响应机制，不仅保障了多模型高并发运行时的响应效率与鲁棒性，还在前端实现了异构推理结果的可视化同屏对比，直观呈现不同模型的中间推理步骤、解析图元、运行耗时及定理调用细节。该系统的成功开发，不仅为多模态几何推理算法提供了统一的测试与分析平台，也有效验证了本文理论方法在实际辅助教学环境中的工程可行性与落地价值。

\xsection{论文组织架构}{Thesis Organization Structure}
本论文的整体内容共划分为六个独立章节，具体的组织安排如下。

第一章为绪论。本章阐述了平面几何求解的研究背景以及在教育方面的研究意义，梳理了几何问题相关数据基准和求解算法的发展脉络，明确了当前研究面临的核心技术挑战，并对本文的具体研究内容进行了详细概述。

第二章介绍了任务概述与相关基础理论。本章对长链条几何推理任务进行了形式化界定，并系统回顾了大模型基本架构、微调范式以及强化学习机制等支撑本文工作的前沿理论。特别是对参数高效微调方法与组相对策略优化算法（GRPO）的深入剖析，为后文两阶段模型训练方案的设计提供了坚实的理论依据。

第三章详细探讨了多模态平面几何基准的构建过程以及基线评估结果。本章阐明了原始试题的采集清洗流程与图文形式化转换方法，展示了数据集的内在统计特征，并建立了标准化的步骤评估协议以完成各类基线模型的性能摸底。通过引入涵盖最终答案准确率与过程逻辑正确率的双重评估标准，本基准揭示了现有主流模型在应对超长推理链条时的能力瓶颈。

第四章提出了面向长链条推理的两阶段增强方法。本章重点阐述了纯文本逻辑增强与多模态过程强化的分离训练思想，阐明了弱符号验证器在提供步骤级反馈中的具体作用，并给出了详实的对比实验与消融分析。实验结果充分显式，该结合强化学习与规则约束的增强方法能够有效缓解视觉特征干扰，显著提升模型多步推演的连贯性与解题成功率。

第五章展示了平面几何问题求解系统的工程设计与应用实现。本章从实际业务需求出发，说明了软件体系的总体架构与调度控制逻辑，并对前后端交互流程及原型系统的功能测试进行了完整展示。依托前后端分离架构与异步并发调度机制，该系统成功实现了多种异构求解模型的高效集成与可视化同屏对比，验证了本文算法的工程落地可行性。

第六章为结论与展望。本章系统总结了本文在多模态几何长链条推理基准构建、求解方法设计与系统实现等方面的研究成果，客观分析了当前方法在复杂空间关系建模、隐性拓扑约束利用等方面的局限，并对几何智能推理领域的后续研究方向进行了展望，为多模态神经符号求解框架的进一步完善提供了参考。
